\documentclass[11pt,a4paper]{article}
\usepackage[margin=18mm]{geometry}
\usepackage{fontspec}
\setmainfont{DejaVu Serif}
\setsansfont{DejaVu Sans}
\setmonofont{DejaVu Sans Mono}
\usepackage{polyglossia}
\setdefaultlanguage{russian}
\usepackage{array,booktabs,tabularx,longtable}
\usepackage{xcolor}
\usepackage{enumitem}
\usepackage{hyperref}
\usepackage{fancyhdr}
\usepackage{titlesec}
\usepackage{listings}
\usepackage{graphicx}
\usepackage{float}
\usepackage{microtype}
\usepackage{parskip}

\definecolor{AltairBlue}{HTML}{1C5FD4}
\definecolor{AltairDark}{HTML}{0B1B2C}
\hypersetup{colorlinks=true,linkcolor=AltairBlue,urlcolor=AltairBlue,pdftitle={ЦУП Альтаир v8 - руководство пользователя}}
\pagestyle{fancy}
\fancyhf{}
\lhead{ЦУП Альтаир v8}
\rhead{Миссия на Луну}
\cfoot{\thepage}
\titleformat{\section}{\Large\bfseries\sffamily\color{AltairBlue}}{\thesection.}{0.5em}{}
\titleformat{\subsection}{\large\bfseries\sffamily}{\thesubsection.}{0.5em}{}
\lstset{basicstyle=\ttfamily\small,breaklines=true,frame=single,columns=fullflexible,showstringspaces=false}
\newcommand{\code}[1]{\texttt{\detokenize{#1}}}

\begin{document}

\begin{center}
{\Huge\bfseries\sffamily ЦУП Альтаир v8}\\[2mm]
{\large Центр управления полётами}\\[3mm]
{\Large\bfseries Миссия на Луну}\\[6mm]
{\large Руководство пользователя и радиошлюза ESP32 + CC1101}\\[4mm]
Версия 0.8.0
\end{center}

\vspace{3mm}
\hrule
\vspace{5mm}

\textbf{Автор: Антон Т.}

ЦУП Альтаир - учебная наземная станция для приёма телеметрии нескольких CubeSat, отображения параметров, построения графиков, просмотра сырых данных последовательного порта и передачи команд через радиошлюз ESP32 + CC1101.

\section{Состав системы}

Типовая цепочка:

\begin{center}
\fbox{\parbox{0.92\linewidth}{\centering
\textbf{CubeSat + CC1101} $\longrightarrow$ \textbf{435.000 МГц} $\longrightarrow$
\textbf{CC1101 + ESP32} $\longrightarrow$ \textbf{USB 115200} $\longrightarrow$
\textbf{ЦУП Альтаир v8}
}}
\end{center}

На площадке Технопрома-2026 система использовалась с крупной наземной антенной, большим экраном ЦУПа и учебными аппаратами с раскрывающимися ленточными антеннами.

\section{Что необходимо скачать}

Для обычного пользователя достаточно:

\begin{enumerate}[leftmargin=7mm]
\item \code{CUP-Altair-Setup-0.8.0.exe} - приложение для Windows.
\item \code{Altair_Gateway_v8.ino} - прошивка ESP32 радиошлюза, если требуется ручная загрузка через Arduino IDE.
\end{enumerate}

Все актуальные файлы находятся в GitHub Release \code{v0.8.0}. Предыдущие GitHub Releases удаляются автоматически при публикации v8, чтобы пользователь видел одну актуальную версию.

\section{Подключение CC1101 к ESP32}

\begin{center}
\begin{tabular}{lll}
\toprule
\textbf{CC1101} & \textbf{ESP32-WROOM-32} & \textbf{Назначение}\\
\midrule
VCC & 3V3 & питание 3.3 В\\
GND & GND & общий провод\\
SCK & GPIO18 & SPI clock\\
MISO / SO & GPIO19 & SPI MISO\\
MOSI / SI & GPIO23 & SPI MOSI\\
CSN / CS & GPIO21 & chip select\\
GDO0 & GPIO4 & пакетный интерфейс\\
GDO2 & GPIO27 & служебный вывод\\
\bottomrule
\end{tabular}
\end{center}

\textbf{Важно:} CC1101 должен питаться от 3.3 В.

\section{Рабочий радиопрофиль}

\begin{center}
\begin{tabular}{ll}
\toprule
Параметр & Значение\\
\midrule
Частота & \textbf{435.000 МГц}\\
Модуляция & 2-FSK\\
Скорость & 4.8 kbps\\
Девиация & 5 кГц\\
Полоса RX & \textbf{203 кГц}\\
Sync word & 0x12AD\\
Encoding & NRZ\\
Длина пакета & variable\\
CRC CC1101 & включён\\
Преамбула & 16 bit\\
USB & 115200 бод\\
\bottomrule
\end{tabular}
\end{center}

Полоса 203 кГц выбрана по результатам реальной работы стенда: часть передатчиков имела заметный частотный уход, и узкая полоса приводила к пропускам.

\section{Формат телеметрии}

Базовый пакет:

\begin{lstlisting}
02,00001,00015,3.00,4.20,1,33
\end{lstlisting}

Поля:

\begin{lstlisting}
ID,PACKET,UPTIME,PANEL_POWER,VOLT,MODE,CHECKSUM
\end{lstlisting}

Расширенный вариант с состоянием антенны:

\begin{lstlisting}
02,00001,00015,3.00,4.20,1,1,33
\end{lstlisting}

Поля:

\begin{lstlisting}
ID,PACKET,UPTIME,PANEL_POWER,VOLT,MODE,ANTENNA,CHECKSUM
\end{lstlisting}

\code{ANTENNA=1} означает раскрытую антенну, \code{ANTENNA=0} - сложенную.

\subsection{Контрольная сумма}

В v8 прикладная проверка XOR \textbf{отключена}. Поле \code{CHECKSUM} принимается и отображается, но пакет не отбрасывается при несовпадении.

Это позволяет участникам менять ID, напряжение, режим и другие параметры без обязательного пересчёта последних двух HEX-символов.

Аппаратный CRC пакетного механизма CC1101 остаётся включённым. Поэтому грубо повреждённый радиопакет по-прежнему может быть отклонён радиочипом.

\section{Подключение к программе}

\begin{enumerate}[leftmargin=7mm]
\item Подключить прошитую ESP32 к компьютеру.
\item Закрыть Arduino Serial Monitor и другие программы, занявшие COM-порт.
\item Запустить ЦУП Альтаир v8.
\item Выбрать профиль ESP32 и скорость 115200 бод.
\item Нажать «Подключить устройство».
\item Выбрать нужный COM-порт.
\end{enumerate}

После успешного подключения панель соединения автоматически сворачивается. Её можно в любой момент развернуть кнопкой в заголовке.

На верхней панели постоянно видна рабочая частота \textbf{435.000 МГц} и полоса приёма 203 кГц.

\section{Несколько спутников}

Графики v8 группируют телеметрию по полю \code{ID}. Каждый новый ID получает собственную линию и цвет.

Например, при одновременном приёме аппаратов 01, 02 и 07 на каждом графике отображаются три отдельные линии. Это устраняет смешивание данных разных спутников.

Кнопка «Сбросить принятые данные» очищает:

\begin{itemize}[leftmargin=7mm]
\item карточки телеметрии;
\item графики и список обнаруженных ID;
\item счётчики пакетов и ошибок;
\item сырой COM-журнал;
\item внутренний буфер парсера.
\end{itemize}

\section{Сырые данные COM}

Панель «Сырые данные последовательного порта» показывает строки, поступающие от ESP32, до их интерпретации интерфейсом.

Она нужна для ответа на главный диагностический вопрос: \textit{приходят ли данные с нужной скоростью вообще?}

Если строки быстро появляются в COM-панели, но карточки не обновляются, проблема находится в формате пакета. Если строк нет, следует проверять радиоканал, ESP32, выбранный порт и питание.

\section{Передача команд}

ESP32-шлюз v8 является приёмопередатчиком. Программа формирует USB-команду:

\begin{lstlisting}
$CMD,RF,TO=02,NAME=PING
\end{lstlisting}

Шлюз передаёт в эфир:

\begin{lstlisting}
$CMD,TO=02,NAME=PING
\end{lstlisting}

Адрес \code{ALL} используется для широковещательной передачи.

\begin{center}
\begin{tabularx}{\linewidth}{lX}
\toprule
\textbf{NAME} & \textbf{Назначение}\\
\midrule
PING & проверка радиопередачи\\
INFO & запрос информации\\
TM\_START & запуск телеметрии\\
TM\_STOP & остановка телеметрии\\
TM\_PERIOD & изменение периода передачи, параметр \code{MS}\\
USER & произвольная команда, параметр \code{DATA}\\
\bottomrule
\end{tabularx}
\end{center}

Бортовая прошивка спутника должна самостоятельно реализовать реакцию на команды. Даже без такой обработки сам факт передачи можно проверить на SDR.

\section{Почему телеметрия может приходить редко}

ЦУП не выполняет искусственного ограничения частоты обновления: разобранная строка сразу передаётся интерфейсу и графикам.

Типовые причины редкого поступления:

\begin{itemize}[leftmargin=7mm]
\item большой период передачи телеметрии на конкретном аппарате;
\item коллизии нескольких передатчиков на одной частоте;
\item частотный уход радиомодуля;
\item ошибки аппаратного CRC CC1101;
\item неправильный COM-порт или занятый Serial Monitor.
\end{itemize}

Для диагностики необходимо сравнить реальный поток в панели COM с ожидаемым периодом передачи.

\section{Первый тест}

После прошивки ESP32 рекомендуется выполнить минимальный тест:

\begin{enumerate}[leftmargin=7mm]
\item Убедиться, что в COM-панели появляется \code{\$INFO,RADIO_READY}.
\item Включить один передатчик.
\item Убедиться, что появляются строки \code{\$TEL,...}.
\item Изменить ID передатчика без пересчёта прикладного XOR: пакет должен продолжить отображаться.
\item Включить второй ID: на графике должна появиться вторая линия.
\item Нажать PING в разделе команд и проверить RF-передачу на SDR или совместимом приёмнике.
\end{enumerate}

\section{О программе}

\textbf{ЦУП Альтаир v8 / 0.8.0.}

Автор: \textbf{Антон Т.}

Программа сделана в рамках проекта «Миссия на Луну» к профильной смене по спутникостроению и развивалась во время Технопрома-2026.

\vfill
\begin{center}
\small ЦУП Альтаир v8 - Миссия на Луну
\end{center}

\end{document}
